% Общий список литературы (подключается через % Общий список литературы (подключается через % Общий список литературы (подключается через % Общий список литературы (подключается через \input{references})
% Нумерация совпадает с literature.md
\begin{thebibliography}{99}
\selectlanguage{english}

\bibitem{elhage2021framework}
N.~Elhage et al. A Mathematical Framework for Transformer Circuits.
Transformer Circuits Thread, 2021.
\url{https://transformer-circuits.pub/2021/framework/index.html}

\bibitem{vaswani2017attention}
A.~Vaswani et al. Attention Is All You Need. NeurIPS, 2017.
arXiv:1706.03762.

\bibitem{geva2021ffn}
M.~Geva et al. Transformer Feed-Forward Layers Are Key-Value Memories.
EMNLP, 2021. arXiv:2012.14913.

\bibitem{cordonnier2020relationship}
J.-B.~Cordonnier et al. On the Relationship between Self-Attention and
Convolutional Layers. ICLR, 2020. arXiv:1911.03584.

\bibitem{dong2021attention}
Y.~Dong et al. Attention is Not All You Need: Pure Attention Loses Rank
Doubly Exponentially with Depth. ICML, 2021. arXiv:2103.03404.

\bibitem{cammarata_clean}
N.~Cammarata. Neural Networks Are Already Clean. Goodfire, 2026.
\url{https://already-clean.goodfire.com/}

\bibitem{gao_notclean}
L.~Gao. Neural networks are not already clean. LessWrong, 2026.
\url{https://www.lesswrong.com/posts/Dkh7pXKjoNWkkxQDz/neural-networks-are-not-already-clean}

\bibitem{olah2020zoom}
C.~Olah et al. Zoom In: An Introduction to Circuits. Distill, 2020.
\url{https://distill.pub/2020/circuits/zoom-in/}

\bibitem{lissack_clean}
M.~Lissack. Already Clean, Not Yet Meaningful: One ``how'' for the
question Cammarata has opened. PhilArchive.
\url{https://philarchive.org/rec/LISACN} % TODO: год

\bibitem{ibm_clean}
S.~Brodsky. Is AI easier to understand than scientists thought? IBM Think, 2026.
\url{https://www.ibm.com/think/news/is-ai-easier-to-understand-than-scientists-thought-goodfire-essay}

\bibitem{katharopoulos2020rnn}
A.~Katharopoulos et al. Transformers are RNNs: Fast Autoregressive
Transformers with Linear Attention. ICML, 2020. arXiv:2006.16236.

\bibitem{wang2020linformer}
S.~Wang et al. Linformer: Self-Attention with Linear Complexity, 2020.
arXiv:2006.04768.

\bibitem{xiong2021nystrom}
Y.~Xiong et al. Nyströmformer: A Nyström-Based Algorithm for
Approximating Self-Attention. AAAI, 2021. arXiv:2102.03902.

\bibitem{choromanski2021performer}
K.~Choromanski et al. Rethinking Attention with Performers. ICLR, 2021.
arXiv:2009.14794.

\bibitem{peng2021rfa}
H.~Peng et al. Random Feature Attention. ICLR, 2021. arXiv:2103.02143.

\bibitem{lu2021soft}
J.~Lu et al. SOFT: Softmax-free Transformer with Linear Complexity.
NeurIPS, 2021. arXiv:2110.11945.

\bibitem{child2019sparse}
R.~Child et al. Generating Long Sequences with Sparse Transformers, 2019.
arXiv:1904.10509.

\bibitem{kitaev2020reformer}
N.~Kitaev et al. Reformer: The Efficient Transformer. ICLR, 2020.
arXiv:2001.04451.

\bibitem{zaheer2020bigbird}
M.~Zaheer et al. Big Bird: Transformers for Longer Sequences. NeurIPS,
2020. arXiv:2007.14062.

\bibitem{leethorp2022fnet}
J.~Lee-Thorp et al. FNet: Mixing Tokens with Fourier Transforms.
NAACL, 2022. arXiv:2105.03824.

\bibitem{tay2021synthesizer}
Y.~Tay et al. Synthesizer: Rethinking Self-Attention in Transformer
Models. ICML, 2021. arXiv:2005.00743.

\bibitem{wong2025a3}
J.~T.~H.~Wong et al. A3: an Analytical Low-Rank Approximation Framework
for Attention, 2025. arXiv:2505.12942.

\bibitem{olsson2022induction}
C.~Olsson et al. In-context Learning and Induction Heads. Transformer
Circuits Thread, 2022. arXiv:2209.11895.

\bibitem{bricken2023monosemanticity}
T.~Bricken et al. Towards Monosemanticity: Decomposing Language Models
With Dictionary Learning. Transformer Circuits Thread, 2023.
\url{https://transformer-circuits.pub/2023/monosemantic-features/index.html}

\bibitem{rai2024review}
D.~Rai et al. A Practical Review of Mechanistic Interpretability for
Transformer-Based Language Models, 2024. arXiv:2407.02646.

\bibitem{yuan2023asvd}
Z.~Yuan et al. ASVD: Activation-aware Singular Value Decomposition for
Compressing Large Language Models, 2023. arXiv:2312.05821.

\bibitem{wang2025svdllm}
X.~Wang et al. SVD-LLM: Truncation-aware Singular Value Decomposition for
Large Language Model Compression. ICLR, 2025. arXiv:2403.07378.

\bibitem{liu2023dejavu}
Z.~Liu et al. Deja Vu: Contextual Sparsity for Efficient LLMs at Inference
Time. ICML, 2023. arXiv:2310.17157.

\bibitem{zhang2022moefication}
Z.~Zhang et al. MoEfication: Transformer Feed-forward Layers are Mixtures of
Experts. Findings of ACL, 2022. arXiv:2110.01786.

\bibitem{lee2024cats}
J.-Y.~Lee et al. CATS: Contextually-Aware Thresholding for Sparsity in Large
Language Models. COLM, 2024. arXiv:2404.08763.

\bibitem{liu2024teal}
J.~Liu et al. Training-Free Activation Sparsity in Large Language Models.
ICLR, 2025. arXiv:2408.14690.

\bibitem{shopkhoev2025replaceme}
D.~Shopkhoev et al. ReplaceMe: Network Simplification via Depth Pruning and
Transformer Block Linearization. NeurIPS, 2025. arXiv:2505.02819.

\bibitem{yun2026ghosted}
V.-D.~Yun et al. Ghosted Layers: Unconstrained Activation Alignment for
Recovering Layer-Pruned LLMs. ICML, 2026. arXiv:2605.15491.

\bibitem{yu2024superweights}
M.~Yu et al. The Super Weight in Large Language Models. arXiv:2411.07191,
2024.

\bibitem{sun2024massive}
M.~Sun et al. Massive Activations in Large Language Models. COLM, 2024.
arXiv:2402.17762.

\bibitem{merity2017pointer}
S.~Merity, C.~Xiong, J.~Bradbury, R.~Socher. Pointer Sentinel Mixture Models. ICLR, 2017.
arXiv:1609.07843. Набор WikiText: \url{https://huggingface.co/datasets/Salesforce/wikitext}

\bibitem{gao2020pile}
L.~Gao et al. The Pile: An 800GB Dataset of Diverse Text for Language Modeling. 2020.
arXiv:2101.00027. Подвыборка pile-10k: \url{https://huggingface.co/datasets/NeelNanda/pile-10k}

\bibitem{gao2023harness}
L.~Gao et al. A framework for few-shot language model evaluation (lm-evaluation-harness), v0.4. Zenodo, 2023.
\url{https://github.com/EleutherAI/lm-evaluation-harness}

\bibitem{paperno2016lambada}
D.~Paperno et al. The LAMBADA dataset: Word prediction requiring a broad discourse context. ACL, 2016.

\bibitem{zellers2019hellaswag}
R.~Zellers, A.~Holtzman, Y.~Bisk, A.~Farhadi, Y.~Choi. HellaSwag: Can a Machine Really Finish Your Sentence? ACL, 2019.

\bibitem{clark2018arc}
P.~Clark et al. Think you have Solved Question Answering? Try ARC, the AI2 Reasoning Challenge. 2018. arXiv:1803.05457.

\bibitem{bisk2020piqa}
Y.~Bisk, R.~Zellers, R.~Le~Bras, J.~Gao, Y.~Choi. PIQA: Reasoning about Physical Commonsense in Natural Language. AAAI, 2020.

\bibitem{sakaguchi2020winogrande}
K.~Sakaguchi, R.~Le~Bras, C.~Bhagavatula, Y.~Choi. WinoGrande: An Adversarial Winograd Schema Challenge at Scale. AAAI, 2020.

\bibitem{hendrycks2021mmlu}
D.~Hendrycks et al. Measuring Massive Multitask Language Understanding. ICLR, 2021. arXiv:2009.03300.

\end{thebibliography}
)
% Нумерация совпадает с literature.md
\begin{thebibliography}{99}
\selectlanguage{english}

\bibitem{elhage2021framework}
N.~Elhage et al. A Mathematical Framework for Transformer Circuits.
Transformer Circuits Thread, 2021.
\url{https://transformer-circuits.pub/2021/framework/index.html}

\bibitem{vaswani2017attention}
A.~Vaswani et al. Attention Is All You Need. NeurIPS, 2017.
arXiv:1706.03762.

\bibitem{geva2021ffn}
M.~Geva et al. Transformer Feed-Forward Layers Are Key-Value Memories.
EMNLP, 2021. arXiv:2012.14913.

\bibitem{cordonnier2020relationship}
J.-B.~Cordonnier et al. On the Relationship between Self-Attention and
Convolutional Layers. ICLR, 2020. arXiv:1911.03584.

\bibitem{dong2021attention}
Y.~Dong et al. Attention is Not All You Need: Pure Attention Loses Rank
Doubly Exponentially with Depth. ICML, 2021. arXiv:2103.03404.

\bibitem{cammarata_clean}
N.~Cammarata. Neural Networks Are Already Clean. Goodfire, 2026.
\url{https://already-clean.goodfire.com/}

\bibitem{gao_notclean}
L.~Gao. Neural networks are not already clean. LessWrong, 2026.
\url{https://www.lesswrong.com/posts/Dkh7pXKjoNWkkxQDz/neural-networks-are-not-already-clean}

\bibitem{olah2020zoom}
C.~Olah et al. Zoom In: An Introduction to Circuits. Distill, 2020.
\url{https://distill.pub/2020/circuits/zoom-in/}

\bibitem{lissack_clean}
M.~Lissack. Already Clean, Not Yet Meaningful: One ``how'' for the
question Cammarata has opened. PhilArchive.
\url{https://philarchive.org/rec/LISACN} % TODO: год

\bibitem{ibm_clean}
S.~Brodsky. Is AI easier to understand than scientists thought? IBM Think, 2026.
\url{https://www.ibm.com/think/news/is-ai-easier-to-understand-than-scientists-thought-goodfire-essay}

\bibitem{katharopoulos2020rnn}
A.~Katharopoulos et al. Transformers are RNNs: Fast Autoregressive
Transformers with Linear Attention. ICML, 2020. arXiv:2006.16236.

\bibitem{wang2020linformer}
S.~Wang et al. Linformer: Self-Attention with Linear Complexity, 2020.
arXiv:2006.04768.

\bibitem{xiong2021nystrom}
Y.~Xiong et al. Nyströmformer: A Nyström-Based Algorithm for
Approximating Self-Attention. AAAI, 2021. arXiv:2102.03902.

\bibitem{choromanski2021performer}
K.~Choromanski et al. Rethinking Attention with Performers. ICLR, 2021.
arXiv:2009.14794.

\bibitem{peng2021rfa}
H.~Peng et al. Random Feature Attention. ICLR, 2021. arXiv:2103.02143.

\bibitem{lu2021soft}
J.~Lu et al. SOFT: Softmax-free Transformer with Linear Complexity.
NeurIPS, 2021. arXiv:2110.11945.

\bibitem{child2019sparse}
R.~Child et al. Generating Long Sequences with Sparse Transformers, 2019.
arXiv:1904.10509.

\bibitem{kitaev2020reformer}
N.~Kitaev et al. Reformer: The Efficient Transformer. ICLR, 2020.
arXiv:2001.04451.

\bibitem{zaheer2020bigbird}
M.~Zaheer et al. Big Bird: Transformers for Longer Sequences. NeurIPS,
2020. arXiv:2007.14062.

\bibitem{leethorp2022fnet}
J.~Lee-Thorp et al. FNet: Mixing Tokens with Fourier Transforms.
NAACL, 2022. arXiv:2105.03824.

\bibitem{tay2021synthesizer}
Y.~Tay et al. Synthesizer: Rethinking Self-Attention in Transformer
Models. ICML, 2021. arXiv:2005.00743.

\bibitem{wong2025a3}
J.~T.~H.~Wong et al. A3: an Analytical Low-Rank Approximation Framework
for Attention, 2025. arXiv:2505.12942.

\bibitem{olsson2022induction}
C.~Olsson et al. In-context Learning and Induction Heads. Transformer
Circuits Thread, 2022. arXiv:2209.11895.

\bibitem{bricken2023monosemanticity}
T.~Bricken et al. Towards Monosemanticity: Decomposing Language Models
With Dictionary Learning. Transformer Circuits Thread, 2023.
\url{https://transformer-circuits.pub/2023/monosemantic-features/index.html}

\bibitem{rai2024review}
D.~Rai et al. A Practical Review of Mechanistic Interpretability for
Transformer-Based Language Models, 2024. arXiv:2407.02646.

\bibitem{yuan2023asvd}
Z.~Yuan et al. ASVD: Activation-aware Singular Value Decomposition for
Compressing Large Language Models, 2023. arXiv:2312.05821.

\bibitem{wang2025svdllm}
X.~Wang et al. SVD-LLM: Truncation-aware Singular Value Decomposition for
Large Language Model Compression. ICLR, 2025. arXiv:2403.07378.

\bibitem{liu2023dejavu}
Z.~Liu et al. Deja Vu: Contextual Sparsity for Efficient LLMs at Inference
Time. ICML, 2023. arXiv:2310.17157.

\bibitem{zhang2022moefication}
Z.~Zhang et al. MoEfication: Transformer Feed-forward Layers are Mixtures of
Experts. Findings of ACL, 2022. arXiv:2110.01786.

\bibitem{lee2024cats}
J.-Y.~Lee et al. CATS: Contextually-Aware Thresholding for Sparsity in Large
Language Models. COLM, 2024. arXiv:2404.08763.

\bibitem{liu2024teal}
J.~Liu et al. Training-Free Activation Sparsity in Large Language Models.
ICLR, 2025. arXiv:2408.14690.

\bibitem{shopkhoev2025replaceme}
D.~Shopkhoev et al. ReplaceMe: Network Simplification via Depth Pruning and
Transformer Block Linearization. NeurIPS, 2025. arXiv:2505.02819.

\bibitem{yun2026ghosted}
V.-D.~Yun et al. Ghosted Layers: Unconstrained Activation Alignment for
Recovering Layer-Pruned LLMs. ICML, 2026. arXiv:2605.15491.

\bibitem{yu2024superweights}
M.~Yu et al. The Super Weight in Large Language Models. arXiv:2411.07191,
2024.

\bibitem{sun2024massive}
M.~Sun et al. Massive Activations in Large Language Models. COLM, 2024.
arXiv:2402.17762.

\bibitem{merity2017pointer}
S.~Merity, C.~Xiong, J.~Bradbury, R.~Socher. Pointer Sentinel Mixture Models. ICLR, 2017.
arXiv:1609.07843. Набор WikiText: \url{https://huggingface.co/datasets/Salesforce/wikitext}

\bibitem{gao2020pile}
L.~Gao et al. The Pile: An 800GB Dataset of Diverse Text for Language Modeling. 2020.
arXiv:2101.00027. Подвыборка pile-10k: \url{https://huggingface.co/datasets/NeelNanda/pile-10k}

\bibitem{gao2023harness}
L.~Gao et al. A framework for few-shot language model evaluation (lm-evaluation-harness), v0.4. Zenodo, 2023.
\url{https://github.com/EleutherAI/lm-evaluation-harness}

\bibitem{paperno2016lambada}
D.~Paperno et al. The LAMBADA dataset: Word prediction requiring a broad discourse context. ACL, 2016.

\bibitem{zellers2019hellaswag}
R.~Zellers, A.~Holtzman, Y.~Bisk, A.~Farhadi, Y.~Choi. HellaSwag: Can a Machine Really Finish Your Sentence? ACL, 2019.

\bibitem{clark2018arc}
P.~Clark et al. Think you have Solved Question Answering? Try ARC, the AI2 Reasoning Challenge. 2018. arXiv:1803.05457.

\bibitem{bisk2020piqa}
Y.~Bisk, R.~Zellers, R.~Le~Bras, J.~Gao, Y.~Choi. PIQA: Reasoning about Physical Commonsense in Natural Language. AAAI, 2020.

\bibitem{sakaguchi2020winogrande}
K.~Sakaguchi, R.~Le~Bras, C.~Bhagavatula, Y.~Choi. WinoGrande: An Adversarial Winograd Schema Challenge at Scale. AAAI, 2020.

\bibitem{hendrycks2021mmlu}
D.~Hendrycks et al. Measuring Massive Multitask Language Understanding. ICLR, 2021. arXiv:2009.03300.

\end{thebibliography}
)
% Нумерация совпадает с literature.md
\begin{thebibliography}{99}
\selectlanguage{english}

\bibitem{elhage2021framework}
N.~Elhage et al. A Mathematical Framework for Transformer Circuits.
Transformer Circuits Thread, 2021.
\url{https://transformer-circuits.pub/2021/framework/index.html}

\bibitem{vaswani2017attention}
A.~Vaswani et al. Attention Is All You Need. NeurIPS, 2017.
arXiv:1706.03762.

\bibitem{geva2021ffn}
M.~Geva et al. Transformer Feed-Forward Layers Are Key-Value Memories.
EMNLP, 2021. arXiv:2012.14913.

\bibitem{cordonnier2020relationship}
J.-B.~Cordonnier et al. On the Relationship between Self-Attention and
Convolutional Layers. ICLR, 2020. arXiv:1911.03584.

\bibitem{dong2021attention}
Y.~Dong et al. Attention is Not All You Need: Pure Attention Loses Rank
Doubly Exponentially with Depth. ICML, 2021. arXiv:2103.03404.

\bibitem{cammarata_clean}
N.~Cammarata. Neural Networks Are Already Clean. Goodfire, 2026.
\url{https://already-clean.goodfire.com/}

\bibitem{gao_notclean}
L.~Gao. Neural networks are not already clean. LessWrong, 2026.
\url{https://www.lesswrong.com/posts/Dkh7pXKjoNWkkxQDz/neural-networks-are-not-already-clean}

\bibitem{olah2020zoom}
C.~Olah et al. Zoom In: An Introduction to Circuits. Distill, 2020.
\url{https://distill.pub/2020/circuits/zoom-in/}

\bibitem{lissack_clean}
M.~Lissack. Already Clean, Not Yet Meaningful: One ``how'' for the
question Cammarata has opened. PhilArchive.
\url{https://philarchive.org/rec/LISACN} % TODO: год

\bibitem{ibm_clean}
S.~Brodsky. Is AI easier to understand than scientists thought? IBM Think, 2026.
\url{https://www.ibm.com/think/news/is-ai-easier-to-understand-than-scientists-thought-goodfire-essay}

\bibitem{katharopoulos2020rnn}
A.~Katharopoulos et al. Transformers are RNNs: Fast Autoregressive
Transformers with Linear Attention. ICML, 2020. arXiv:2006.16236.

\bibitem{wang2020linformer}
S.~Wang et al. Linformer: Self-Attention with Linear Complexity, 2020.
arXiv:2006.04768.

\bibitem{xiong2021nystrom}
Y.~Xiong et al. Nyströmformer: A Nyström-Based Algorithm for
Approximating Self-Attention. AAAI, 2021. arXiv:2102.03902.

\bibitem{choromanski2021performer}
K.~Choromanski et al. Rethinking Attention with Performers. ICLR, 2021.
arXiv:2009.14794.

\bibitem{peng2021rfa}
H.~Peng et al. Random Feature Attention. ICLR, 2021. arXiv:2103.02143.

\bibitem{lu2021soft}
J.~Lu et al. SOFT: Softmax-free Transformer with Linear Complexity.
NeurIPS, 2021. arXiv:2110.11945.

\bibitem{child2019sparse}
R.~Child et al. Generating Long Sequences with Sparse Transformers, 2019.
arXiv:1904.10509.

\bibitem{kitaev2020reformer}
N.~Kitaev et al. Reformer: The Efficient Transformer. ICLR, 2020.
arXiv:2001.04451.

\bibitem{zaheer2020bigbird}
M.~Zaheer et al. Big Bird: Transformers for Longer Sequences. NeurIPS,
2020. arXiv:2007.14062.

\bibitem{leethorp2022fnet}
J.~Lee-Thorp et al. FNet: Mixing Tokens with Fourier Transforms.
NAACL, 2022. arXiv:2105.03824.

\bibitem{tay2021synthesizer}
Y.~Tay et al. Synthesizer: Rethinking Self-Attention in Transformer
Models. ICML, 2021. arXiv:2005.00743.

\bibitem{wong2025a3}
J.~T.~H.~Wong et al. A3: an Analytical Low-Rank Approximation Framework
for Attention, 2025. arXiv:2505.12942.

\bibitem{olsson2022induction}
C.~Olsson et al. In-context Learning and Induction Heads. Transformer
Circuits Thread, 2022. arXiv:2209.11895.

\bibitem{bricken2023monosemanticity}
T.~Bricken et al. Towards Monosemanticity: Decomposing Language Models
With Dictionary Learning. Transformer Circuits Thread, 2023.
\url{https://transformer-circuits.pub/2023/monosemantic-features/index.html}

\bibitem{rai2024review}
D.~Rai et al. A Practical Review of Mechanistic Interpretability for
Transformer-Based Language Models, 2024. arXiv:2407.02646.

\bibitem{yuan2023asvd}
Z.~Yuan et al. ASVD: Activation-aware Singular Value Decomposition for
Compressing Large Language Models, 2023. arXiv:2312.05821.

\bibitem{wang2025svdllm}
X.~Wang et al. SVD-LLM: Truncation-aware Singular Value Decomposition for
Large Language Model Compression. ICLR, 2025. arXiv:2403.07378.

\bibitem{liu2023dejavu}
Z.~Liu et al. Deja Vu: Contextual Sparsity for Efficient LLMs at Inference
Time. ICML, 2023. arXiv:2310.17157.

\bibitem{zhang2022moefication}
Z.~Zhang et al. MoEfication: Transformer Feed-forward Layers are Mixtures of
Experts. Findings of ACL, 2022. arXiv:2110.01786.

\bibitem{lee2024cats}
J.-Y.~Lee et al. CATS: Contextually-Aware Thresholding for Sparsity in Large
Language Models. COLM, 2024. arXiv:2404.08763.

\bibitem{liu2024teal}
J.~Liu et al. Training-Free Activation Sparsity in Large Language Models.
ICLR, 2025. arXiv:2408.14690.

\bibitem{shopkhoev2025replaceme}
D.~Shopkhoev et al. ReplaceMe: Network Simplification via Depth Pruning and
Transformer Block Linearization. NeurIPS, 2025. arXiv:2505.02819.

\bibitem{yun2026ghosted}
V.-D.~Yun et al. Ghosted Layers: Unconstrained Activation Alignment for
Recovering Layer-Pruned LLMs. ICML, 2026. arXiv:2605.15491.

\bibitem{yu2024superweights}
M.~Yu et al. The Super Weight in Large Language Models. arXiv:2411.07191,
2024.

\bibitem{sun2024massive}
M.~Sun et al. Massive Activations in Large Language Models. COLM, 2024.
arXiv:2402.17762.

\bibitem{merity2017pointer}
S.~Merity, C.~Xiong, J.~Bradbury, R.~Socher. Pointer Sentinel Mixture Models. ICLR, 2017.
arXiv:1609.07843. Набор WikiText: \url{https://huggingface.co/datasets/Salesforce/wikitext}

\bibitem{gao2020pile}
L.~Gao et al. The Pile: An 800GB Dataset of Diverse Text for Language Modeling. 2020.
arXiv:2101.00027. Подвыборка pile-10k: \url{https://huggingface.co/datasets/NeelNanda/pile-10k}

\bibitem{gao2023harness}
L.~Gao et al. A framework for few-shot language model evaluation (lm-evaluation-harness), v0.4. Zenodo, 2023.
\url{https://github.com/EleutherAI/lm-evaluation-harness}

\bibitem{paperno2016lambada}
D.~Paperno et al. The LAMBADA dataset: Word prediction requiring a broad discourse context. ACL, 2016.

\bibitem{zellers2019hellaswag}
R.~Zellers, A.~Holtzman, Y.~Bisk, A.~Farhadi, Y.~Choi. HellaSwag: Can a Machine Really Finish Your Sentence? ACL, 2019.

\bibitem{clark2018arc}
P.~Clark et al. Think you have Solved Question Answering? Try ARC, the AI2 Reasoning Challenge. 2018. arXiv:1803.05457.

\bibitem{bisk2020piqa}
Y.~Bisk, R.~Zellers, R.~Le~Bras, J.~Gao, Y.~Choi. PIQA: Reasoning about Physical Commonsense in Natural Language. AAAI, 2020.

\bibitem{sakaguchi2020winogrande}
K.~Sakaguchi, R.~Le~Bras, C.~Bhagavatula, Y.~Choi. WinoGrande: An Adversarial Winograd Schema Challenge at Scale. AAAI, 2020.

\bibitem{hendrycks2021mmlu}
D.~Hendrycks et al. Measuring Massive Multitask Language Understanding. ICLR, 2021. arXiv:2009.03300.

\end{thebibliography}
)
% Нумерация совпадает с literature.md
\begin{thebibliography}{99}
\selectlanguage{english}

\bibitem{elhage2021framework}
N.~Elhage et al. A Mathematical Framework for Transformer Circuits.
Transformer Circuits Thread, 2021.
\url{https://transformer-circuits.pub/2021/framework/index.html}

\bibitem{vaswani2017attention}
A.~Vaswani et al. Attention Is All You Need. NeurIPS, 2017.
arXiv:1706.03762.

\bibitem{geva2021ffn}
M.~Geva et al. Transformer Feed-Forward Layers Are Key-Value Memories.
EMNLP, 2021. arXiv:2012.14913.

\bibitem{cordonnier2020relationship}
J.-B.~Cordonnier et al. On the Relationship between Self-Attention and
Convolutional Layers. ICLR, 2020. arXiv:1911.03584.

\bibitem{dong2021attention}
Y.~Dong et al. Attention is Not All You Need: Pure Attention Loses Rank
Doubly Exponentially with Depth. ICML, 2021. arXiv:2103.03404.

\bibitem{cammarata_clean}
N.~Cammarata. Neural Networks Are Already Clean. Goodfire, 2026.
\url{https://already-clean.goodfire.com/}

\bibitem{gao_notclean}
L.~Gao. Neural networks are not already clean. LessWrong, 2026.
\url{https://www.lesswrong.com/posts/Dkh7pXKjoNWkkxQDz/neural-networks-are-not-already-clean}

\bibitem{olah2020zoom}
C.~Olah et al. Zoom In: An Introduction to Circuits. Distill, 2020.
\url{https://distill.pub/2020/circuits/zoom-in/}

\bibitem{lissack_clean}
M.~Lissack. Already Clean, Not Yet Meaningful: One ``how'' for the
question Cammarata has opened. PhilArchive.
\url{https://philarchive.org/rec/LISACN} % TODO: год

\bibitem{ibm_clean}
S.~Brodsky. Is AI easier to understand than scientists thought? IBM Think, 2026.
\url{https://www.ibm.com/think/news/is-ai-easier-to-understand-than-scientists-thought-goodfire-essay}

\bibitem{katharopoulos2020rnn}
A.~Katharopoulos et al. Transformers are RNNs: Fast Autoregressive
Transformers with Linear Attention. ICML, 2020. arXiv:2006.16236.

\bibitem{wang2020linformer}
S.~Wang et al. Linformer: Self-Attention with Linear Complexity, 2020.
arXiv:2006.04768.

\bibitem{xiong2021nystrom}
Y.~Xiong et al. Nyströmformer: A Nyström-Based Algorithm for
Approximating Self-Attention. AAAI, 2021. arXiv:2102.03902.

\bibitem{choromanski2021performer}
K.~Choromanski et al. Rethinking Attention with Performers. ICLR, 2021.
arXiv:2009.14794.

\bibitem{peng2021rfa}
H.~Peng et al. Random Feature Attention. ICLR, 2021. arXiv:2103.02143.

\bibitem{lu2021soft}
J.~Lu et al. SOFT: Softmax-free Transformer with Linear Complexity.
NeurIPS, 2021. arXiv:2110.11945.

\bibitem{child2019sparse}
R.~Child et al. Generating Long Sequences with Sparse Transformers, 2019.
arXiv:1904.10509.

\bibitem{kitaev2020reformer}
N.~Kitaev et al. Reformer: The Efficient Transformer. ICLR, 2020.
arXiv:2001.04451.

\bibitem{zaheer2020bigbird}
M.~Zaheer et al. Big Bird: Transformers for Longer Sequences. NeurIPS,
2020. arXiv:2007.14062.

\bibitem{leethorp2022fnet}
J.~Lee-Thorp et al. FNet: Mixing Tokens with Fourier Transforms.
NAACL, 2022. arXiv:2105.03824.

\bibitem{tay2021synthesizer}
Y.~Tay et al. Synthesizer: Rethinking Self-Attention in Transformer
Models. ICML, 2021. arXiv:2005.00743.

\bibitem{wong2025a3}
J.~T.~H.~Wong et al. A3: an Analytical Low-Rank Approximation Framework
for Attention, 2025. arXiv:2505.12942.

\bibitem{olsson2022induction}
C.~Olsson et al. In-context Learning and Induction Heads. Transformer
Circuits Thread, 2022. arXiv:2209.11895.

\bibitem{bricken2023monosemanticity}
T.~Bricken et al. Towards Monosemanticity: Decomposing Language Models
With Dictionary Learning. Transformer Circuits Thread, 2023.
\url{https://transformer-circuits.pub/2023/monosemantic-features/index.html}

\bibitem{rai2024review}
D.~Rai et al. A Practical Review of Mechanistic Interpretability for
Transformer-Based Language Models, 2024. arXiv:2407.02646.

\bibitem{yuan2023asvd}
Z.~Yuan et al. ASVD: Activation-aware Singular Value Decomposition for
Compressing Large Language Models, 2023. arXiv:2312.05821.

\bibitem{wang2025svdllm}
X.~Wang et al. SVD-LLM: Truncation-aware Singular Value Decomposition for
Large Language Model Compression. ICLR, 2025. arXiv:2403.07378.

\bibitem{liu2023dejavu}
Z.~Liu et al. Deja Vu: Contextual Sparsity for Efficient LLMs at Inference
Time. ICML, 2023. arXiv:2310.17157.

\bibitem{zhang2022moefication}
Z.~Zhang et al. MoEfication: Transformer Feed-forward Layers are Mixtures of
Experts. Findings of ACL, 2022. arXiv:2110.01786.

\bibitem{lee2024cats}
J.-Y.~Lee et al. CATS: Contextually-Aware Thresholding for Sparsity in Large
Language Models. COLM, 2024. arXiv:2404.08763.

\bibitem{liu2024teal}
J.~Liu et al. Training-Free Activation Sparsity in Large Language Models.
ICLR, 2025. arXiv:2408.14690.

\bibitem{shopkhoev2025replaceme}
D.~Shopkhoev et al. ReplaceMe: Network Simplification via Depth Pruning and
Transformer Block Linearization. NeurIPS, 2025. arXiv:2505.02819.

\bibitem{yun2026ghosted}
V.-D.~Yun et al. Ghosted Layers: Unconstrained Activation Alignment for
Recovering Layer-Pruned LLMs. ICML, 2026. arXiv:2605.15491.

\bibitem{yu2024superweights}
M.~Yu et al. The Super Weight in Large Language Models. arXiv:2411.07191,
2024.

\bibitem{sun2024massive}
M.~Sun et al. Massive Activations in Large Language Models. COLM, 2024.
arXiv:2402.17762.

\bibitem{merity2017pointer}
S.~Merity, C.~Xiong, J.~Bradbury, R.~Socher. Pointer Sentinel Mixture Models. ICLR, 2017.
arXiv:1609.07843. Набор WikiText: \url{https://huggingface.co/datasets/Salesforce/wikitext}

\bibitem{gao2020pile}
L.~Gao et al. The Pile: An 800GB Dataset of Diverse Text for Language Modeling. 2020.
arXiv:2101.00027. Подвыборка pile-10k: \url{https://huggingface.co/datasets/NeelNanda/pile-10k}

\bibitem{gao2023harness}
L.~Gao et al. A framework for few-shot language model evaluation (lm-evaluation-harness), v0.4. Zenodo, 2023.
\url{https://github.com/EleutherAI/lm-evaluation-harness}

\bibitem{paperno2016lambada}
D.~Paperno et al. The LAMBADA dataset: Word prediction requiring a broad discourse context. ACL, 2016.

\bibitem{zellers2019hellaswag}
R.~Zellers, A.~Holtzman, Y.~Bisk, A.~Farhadi, Y.~Choi. HellaSwag: Can a Machine Really Finish Your Sentence? ACL, 2019.

\bibitem{clark2018arc}
P.~Clark et al. Think you have Solved Question Answering? Try ARC, the AI2 Reasoning Challenge. 2018. arXiv:1803.05457.

\bibitem{bisk2020piqa}
Y.~Bisk, R.~Zellers, R.~Le~Bras, J.~Gao, Y.~Choi. PIQA: Reasoning about Physical Commonsense in Natural Language. AAAI, 2020.

\bibitem{sakaguchi2020winogrande}
K.~Sakaguchi, R.~Le~Bras, C.~Bhagavatula, Y.~Choi. WinoGrande: An Adversarial Winograd Schema Challenge at Scale. AAAI, 2020.

\bibitem{hendrycks2021mmlu}
D.~Hendrycks et al. Measuring Massive Multitask Language Understanding. ICLR, 2021. arXiv:2009.03300.

\end{thebibliography}
